% Generated draft: requires mathematical and source review.
\documentclass[11pt]{article}
\usepackage[margin=25mm]{geometry}
\usepackage{fontspec}
\setmainfont{Latin Modern Roman}
\usepackage{amsmath,amssymb,mathtools,unicode-math}
\setmathfont{Latin Modern Math}
\usepackage{enumitem,needspace,xurl,hyperref}
\setlength{\parindent}{0pt}
\setlength{\parskip}{5pt}
\setlength{\emergencystretch}{3em}
\newcommand{\N}{\mathbb{N}}
\newcommand{\Z}{\mathbb{Z}}
\newcommand{\Q}{\mathbb{Q}}
\newcommand{\R}{\mathbb{R}}
\newcommand{\C}{\mathbb{C}}
\newcommand{\F}{\mathbb{F}}
\newcommand{\A}{\mathbb{A}}
\newcommand{\PP}{\mathbb{P}}
\newcommand{\E}{\mathbb{E}}
\newcommand{\eps}{\varepsilon}
\newcommand{\dd}{\,\mathrm{d}}
\begin{document}
\section*{UnsolvedMath ID: 30004098}
\textbf{Source problem code:} catalogue entry
\subsection{Definitions and mathematical statement}
The supplied question asks whether the finite-type property of finite tensor categories is preserved under duality. The record does not specify the ground field, the definition of finite type, or the intended duality. The following is a conventional formulation requiring confirmation against the cited sources.

Fix an algebraically closed field \(k\). A finite tensor category \(\mathcal C\) over \(k\) is a rigid \(k\)-linear abelian monoidal category whose underlying abelian category is equivalent to the category of finite-dimensional modules over a finite-dimensional \(k\)-algebra, whose tensor product is \(k\)-bilinear and exact in each variable, and whose tensor unit \(\mathbf 1\) is simple. Set
\[
H^\bullet(\mathcal C,V)=\bigoplus_{n\geq 0}\operatorname{Ext}^n_{\mathcal C}(\mathbf 1,V),
\qquad
H^\bullet(\mathcal C)=H^\bullet(\mathcal C,\mathbf 1).
\]
The latter is an algebra under Yoneda composition, and the former is a graded module over it by Yoneda composition. In the cohomological interpretation used here, \(\mathcal C\) is of finite type if \(H^\bullet(\mathcal C)\) is finitely generated as a \(k\)-algebra and \(H^\bullet(\mathcal C,V)\) is finitely generated over \(H^\bullet(\mathcal C)\) for every object \(V\in\mathcal C\).

A finite left \(\mathcal C\)-module category \(\mathcal M\) is a finite \(k\)-linear abelian category equipped with a compatible action of \(\mathcal C\). It is exact if \(P\otimes M\) is projective in \(\mathcal M\) for every projective object \(P\in\mathcal C\) and every \(M\in\mathcal M\). It is indecomposable if it is not a direct sum of two nonzero \(\mathcal C\)-module categories. For a nonzero exact indecomposable \(\mathcal M\), write
\[
\mathcal C_{\mathcal M}^{*}=\operatorname{Fun}_{\mathcal C}^{\mathrm{rex}}(\mathcal M,\mathcal M)
\]
for the dual tensor category of right exact \(\mathcal C\)-module endofunctors, with tensor product given by composition and unit the identity functor.

\textbf{Candidate precise question under these conventions.} For every algebraically closed field \(k\), every finite tensor category \(\mathcal C\) over \(k\), and every nonzero exact indecomposable finite left \(\mathcal C\)-module category \(\mathcal M\), does
\[
\mathcal C\text{ is of finite type}
\quad\Longrightarrow\quad
\mathcal C_{\mathcal M}^{*}\text{ is of finite type}
\]
hold? Here the conclusion includes finite generation with coefficients in every object of \(\mathcal C_{\mathcal M}^{*}\).

The universal field quantifier and the module-category interpretation of duality are explicit reconstruction choices, not hypotheses established by the supplied statement. The intended characteristic restrictions, duality convention, and precise finite-generation requirements remain to be checked.
\subsection{Short English statement}
The question asks whether finite generation of tensor-category cohomology, including cohomology with all object coefficients, survives passage to a dual tensor category. The displayed formulation interprets duality through exact indecomposable module categories; the supplied record does not specify these conventions. Its supplied status is \emph{partially solved}, with an assessment that broad special cases are known while the unrestricted conjecture remains open; this assessment is dated curation, not independently verified here.
\subsection{Sources}
\begin{itemize}
\item Supplied local registry record, permanent UnsolvedMath ID \(30004098\), source problem code \texttt{catalogue entry}, title \emph{Duality Invariance of Finite Type Tensor Categories}; collection: Oberwolfach Reports Open Problems. No individual proposer is supplied; the proposed year is recorded as 2019.
\item Supplied source citation: \emph{Mini-Workshop: Cohomology of Hopf Algebras and Tensor Categories (2020)}, Oberwolfach Reports, pp.\ 663--693, DOI \texttt{10.4171/owr/2019/11}. \url{https://doi.org/10.4171/owr/2019/11}. The report is described as documenting mathematical workshops and their open-problem discussions.
\item Supplied literature reference: \emph{Cohomology and duality for finite tensor categories}. \url{https://doi.org/10.1090/tran/8548}.
\item Supplied literature reference: Negron--Plavnik preprint. \url{https://arxiv.org/abs/1807.08854}.
\item Supplied literature reference: Oberwolfach report. \url{https://doi.org/10.4171/owr/2019/11}.
\item Literature assessment dated exactly \texttt{2026-08-22}: the supplied curation labels the problem partially solved and reports preservation of finite-type cohomology in several broad settings, including suitable exact module categories and important braided examples, while describing the unrestricted finite-tensor-category conjecture as open. This is supplied dated curation derived from the OWR record, not an independently verified assessment; its time-sensitive status and cited claims require independent verification.
\item The registry labels the statement \texttt{corrected\_verified} and says its mathematical content was preserved while notation, typography, and sentence structure were repaired for publication. This is supplied editorial metadata, not mathematical verification of this reconstructed formulation.
\end{itemize}
\end{document}
